\begin{exercise}{50}
Show that the composite
\[
 \mathrm{Cat}_{\Delta}\xrightarrow{\ \mathrm{N}\ }\mathrm{sSet}
 \xrightarrow{\ h\ }\mathrm{Cat}
\]
is isomorphic to the functor $\pi\colon\mathrm{Cat}_{\Delta}\to\mathrm{Cat}$.

\begin{proof}
The vertices of $\mathrm{N}(\mathcal C)$ are the objects of
$\mathcal C$, its edges $x\to y$ are the vertices of
$\mathrm{Hom}_{\mathcal C}(x,y)$, and degenerate edges are units.
By Definition~1.2.56, a triangle with successive edges $f,g$
and long edge $k$ is specified by a hom-edge $k\to gf$.
Thus the presentation of $h$ in Definition~1.2.5 gives a functor
\[
 \Phi_{\mathcal C}\colon h\mathrm{N}(\mathcal C)\longrightarrow\pi\mathcal C
\]
which is the identity on objects and sends the class of an edge $f$
to its connected component $[f]$.
Indeed, degenerate edges map to identities, and every triangle
gives $[k]=[gf]=[g][f]$ in $\pi\mathcal C$
(Definition~1.2.48).

Conversely, send a component $[f]$ to the edge class
$\langle f\rangle$ in $h\mathrm{N}(\mathcal C)$.
For a hom-edge $f\to g$, the triangle from Exercise~\ref{ex:42}
has successive edges $\mathrm{id}_x,g$ and long edge $f$;
hence $\langle f\rangle=\langle g\rangle$.
Since connected components are generated by zigzags of hom-edges,
this assignment is well defined.
The constant hom-edge at $gf$ gives a triangle with successive
edges $f,g$ and long edge $gf$, so
$\langle gf\rangle=\langle g\rangle\langle f\rangle$.
Together with the degenerate-edge relation, this shows that
the assignment defines a functor
$\Psi_{\mathcal C}\colon\pi\mathcal C\to h\mathrm{N}(\mathcal C)$.

The two composites are identities on objects and on the edge
classes, which generate $h\mathrm{N}(\mathcal C)$, and on every
component $[f]$. Thus $\Phi_{\mathcal C}$ is an isomorphism.
A simplicial functor $F$ sends each edge $f$ to $F(f)$,
so the comparisons commute with $F$ on objects and generators.
They therefore form the required natural isomorphism.
\end{proof}
\end{exercise}
